\section{Related Work}
\label{sec:related_work}

\para{In-context graph representations.}
The closest motivation is the graph tracing setup of \citet{park2025iclr}. Their work studies whether language-model representations reorganize around context-specified square grids, rings, and other graph structures when a prompt supplies random-walk evidence. Our work keeps the same broad object of study, namely novel graph geometry induced inside a context, but asks a different question. Rather than measuring internal activation geometry or next-token continuation, we ask whether the induced relation is verbally addressable through explicit adjacency queries. This behavioral readout is important because a representation can be structured without being reliably exposed by a prompt.

\para{Abstract visual and spatial reasoning.}
Several benchmarks show that current multimodal systems struggle with controlled geometric and spatial structure. Bongard-LOGO defines concept-learning problems over procedurally generated 2D shapes with underlying programs \cite{nie2020bongard}. Shape-Blind shows failures on simple polygon recognition and side counting \cite{rudman2025forgotten}. PuzzleVQA diagnoses abstract visual pattern reasoning by separating perception, induction, and deduction \cite{chia2024puzzlevqa}. Visual Spatial Reasoning tests natural-image spatial relations \cite{liu2023visual}, while VisuLogic and Spatial-DISE broaden the scope to visual-centric logic and spatial transformations \cite{xu2025visulogic,huang2025spatial}. These datasets motivate careful controls for geometry, but most evaluate visual input or multiple-choice classification. We instead use text-only graph traces so that the uncertainty is not image perception but the prompt-accessibility of a newly defined relation.

\para{Prompting, symbolic grounding, and repetition.}
Work on Bongard-style reasoning finds that direct prompting can underperform structured or descriptive inputs \cite{malkinski2024reasoning}. Symbolic grounding work further shows that models can reason better when abstract visual concepts are represented as programs or language descriptions rather than pixels \cite{vaishnav2026symbolic}. These results suggest that failure can come from the interface between representation and query, not only from high-level reasoning. Our study follows the same diagnostic spirit. We compare a terse direct prompt against an explicit-reasoning prompt, but we also add a memorization baseline. This baseline is central: if performance is explained by transitions already observed in context, then the model is addressable as a repeated-edge memory, not as an inferred topology.
